\textbf{Proof.}
Assumption ass:coordinate-bounded-kappa-feasibility ensures that the two pilot-selected bases in def:feasible-ht-pilot-selected-estimator are well defined in the nonempty class \(\mathcal B_{\kappa,\overline a}\).  For the \(n\)-th problem, put
\[
 v_n=\inf_{B\in\mathcal B_Z}w(B),
 \qquad
 q_n=\inf_{B\in\mathcal A_{0,Z}}w(B).
 \tag{1}
\]
By hypothesis, every attained full ordinary-domain minimax and zero-covariance optimizer lies in \(\mathcal B_{\kappa,\overline a}\).  Since \(\mathcal B_{\kappa,\overline a}\subseteq\mathcal B_Z\), a full minimax optimizer is feasible for the restricted problem, and therefore
\[
 v_{\mathrm{mm},\overline a}=v_n.
 \tag{2}
\]
Likewise, a full zero-covariance optimizer is feasible for the restricted maximization, so the restricted and full maxima agree.  Every restricted maximizer then attains the full maximum and is a full maximizer, while the assumed containment places every full maximizer in the restricted class.  Hence
\[
 \mathcal A_{0,\overline a}=\mathcal A_{0,Z}.
 \tag{3}
\]
By def:ordinary-intercept-slope-image and thm:ordinary-inverse-strict-improvement-characterization,
\[
 \Delta_Z=q_n-v_n.
 \tag{4}
\]

Use the notation \(M_n\), \(C_n\), \(L_n\), \(R_n\), and \(V_s\) from the finite inequalities in thm:pilot-selected-ht-first-order-variance-transfer.  Its robust-selector inequality and its zero-surrogate inequality, together with (2)--(4), imply
\[
 \mathbb E_Pw(\widehat B_0)-\mathbb E_Pw(\widehat B_{\mathrm{mm}})
 \geq
 \Delta_Z-(4L_n+2)C_n\mathbb E_PM_n.
 \tag{5}
\]
The two applications of its feasible-estimator bridge give
\[
 V_0\geq\mathbb E_Pw(\widehat B_0)-R_n,
 \qquad
 V_{\mathrm{mm}}\leq\mathbb E_Pw(\widehat B_{\mathrm{mm}})+R_n.
 \tag{6}
\]
Combining (5) and (6) proves
\[
 \boxed{
 V_0-V_{\mathrm{mm}}
 \geq
 \Delta_Z-(4L_n+2)C_n\mathbb E_PM_n-2R_n.}
 \tag{7}
\]
Every term in (7) is either the ordinary-domain population gap, a pilot endpoint sensitivity term, or the explicit feasible-coefficient remainder; no desired variance ordering has been assumed.

We next turn (7) into a directly checkable sufficient condition.  Lemma lem:uniform-pilot-endpoint-matrix-rate supplies finite constants \(K_M,K_C<\infty\), independent of \(n\), such that for all sufficiently large \(n\), with \(m_{\min}=m_1\wedge m_0\),
\[
 \mathbb E_PM_n\leq K_M\frac{n}{\sqrt{m_{\min}}},
 \qquad
 C_n\leq K_Cn^{\alpha+2\beta-1}.
 \tag{8}
\]
Lemma lem:uniform-zhao-ht-oracle-variance-equivalence supplies \(R_n\leq K_Rn^{4\alpha+4\beta-2}\).  Substitution in (7) gives
\[
V_0-V_{\mathrm{mm}}
\ge
\Delta_Z
-(4L_n+2)K_CK_M\frac{n^{\alpha+2\beta}}{\sqrt{m_{\min}}}
-2K_Rn^{4\alpha+4\beta-2}.
\tag{9}
\]
Therefore the displayed strictness condition in the proposition implies \(V_0>V_{\mathrm{mm}}\) at that finite \(n\).  Any certified constants satisfying (8) and the remainder bound may be used, so this is an audit condition rather than an asymptotic relabeling.

Finally, the exact sensitivity condition in the inherited theorem assumptions gives
\[
 nC_n\mathbb E_PM_n\longrightarrow0.
 \tag{10}
\]
The boundedness of \(L_n\) preserves (10) after multiplication by \(4L_n+2\).  Moreover, \(4\alpha+4\beta<1\) and the remainder bound give
\[
 nR_n\leq K_Rn^{4\alpha+4\beta-1}\longrightarrow0.
 \tag{11}
\]
Multiplying (7) by \(n\), taking lower limits, and using (10)--(11) yields
\[
 \liminf_{n\to\infty}n(V_0-V_{\mathrm{mm}})
 \geq\liminf_{n\to\infty}n\Delta_Z>0.
\]
The variances in \(V_0\) and \(V_{\mathrm{mm}}\) remain conditional design variances under \(\nu\); the pilot and coupling expectations are only ex ante evaluations.  No statement about the self-normalized variant, a Wald law, or a replicated-triangle sequence is used. \(\square\)